\chapter{Wstęp}

\label{ch:wstep}


\section{Wprowadzenie}

Dynamiczny rozwój rynku kryptowalut oraz jego wysoka zmienność powodują, że~organizacje finansowe potrzebują wiarygodnych strumieni danych o~niskim opóźnieniu.
Rynek ten generuje dzienny wolumen transakcyjny rzędu dziesiątek miliardów dolarów~\cite{coingecko}.
Jednocześnie specyfika giełd cyfrowych niesie ze~sobą ryzyka związane z~jakością danych.
Zjawiska takie jak \textit{wash trading}, czyli sztuczne generowanie obrotu, mogą prowadzić do~zafałszowania informacji rynkowych, przez co~publicznie dostępne, niesprawdzone dane mogą okazać się niewiarygodne~\cite{wash_trading}.
W~konsekwencji podmioty analizujące rynek, w~szczególności w~kontekście wykrywania anomalii czy~okazji arbitrażowych, nie~mogą opierać się wyłącznie na~zewnętrznych agregatorach, lecz powinny budować własne potoki przetwarzania danych (\textit{data pipelines})~\cite{arbitrage_finance}.

Ciągły dostęp do~danych historycznych i~strumieniowych jest fundamentem algorytmicznych strategii inwestycyjnych.
W~środowiskach o~ograniczonym zaufaniu do~sieci zewnętrznej istotna staje się także możliwość autonomicznego działania systemu (\textit{disconnected operations}), co~wpisuje się w~założenia przetwarzania brzegowego (\textit{edge computing})~\cite{edge_computing_satya}.

Projektowana architektura obejmuje serwer strumieniowy nasłuchujący kanałów WebSocket giełdy Binance, usługę wymiany komunikatów JetStream, moduły \textit{ETL} odpowiedzialne za~transformację i~wzbogacanie danych oraz warstwę analityczną udostępniającą zagregowane metryki.
Dostęp do~danych zapewnia interfejs GraphQL, natomiast bot platformy Telegram odpowiada za~dystrybucję powiadomień o~zdarzeniach rynkowych w~czasie rzeczywistym.

Kluczowym elementem infrastruktury jest mechanizm \textit{tiered storage} (pamięć warstwowa).
Dane w~systemie trafiają najpierw do~warstwy gorącej (MongoDB), a~następnie~-- po~upływie określonego czasu retencji~-- przenoszone są do~warstwy zimnej opartej na~MinIO. Pozwala to~odciążyć bazę operacyjną, obniżyć koszty przechowywania danych historycznych oraz zwiększyć niezawodność przetwarzania~\cite{unity_tiered_storage}.
W~warstwie obiektowej dane mogą być dodatkowo kompresowane i~wersjonowane.


\section{Cel pracy}

Celem pracy jest zaprojektowanie i~implementacja kompletnego łańcucha przetwarzania danych rynkowych obejmującego odbiór, walidację, wzbogacanie, udostępnianie oraz archiwizację informacji w~trybie ciągłym.
W~ramach pracy został opracowany niezawodny serwer strumieniowy utrzymujący połączenia z~WebSocket Binance.
Zaprojektowano kanał dystrybucyjny oparty na~NATS JetStream, zapewniający buforowanie zdarzeń oraz odporność na~awarie usług zależnych.
Zaimplementowano procesy \textit{ETL} dzielące dane na~zbiory surowe i~oczyszczone oraz wzbogacające je~o~metadane analityczne.
Skonfigurowano i~zautomatyzowano przenoszenie danych w~modelu \textit{tiered storage} (MongoDB~→~MinIO), zapewniając optymalizację zasobów warstwy operacyjnej przy równoczesnej dostępności danych historycznych.
Zintegrowano system powiadomień z~Telegramem, umożliwiając przesyłanie alertów rynkowych w~czasie rzeczywistym.

Realizacja założeń obejmuje przygotowanie dokumentacji architektury, implementację usług w~języku \textit{Go} oraz konfigurację infrastruktury w~środowisku Docker.


\section{Plan pracy}

Niniejsza praca składa się z~siedmiu rozdziałów, z~czego pierwszy stanowi wprowadzenie do~problematyki, kolejne cztery opisują realizację systemu, zaś dwa ostatnie podsumowują osiągnięte rezultaty oraz wskazują kierunki dalszego rozwoju.

W~rozdziale~\ref{ch:analiza} przedstawiono analizę wymagań funkcjonalnych i~niefunkcjonalnych systemu, obejmującą m.in.\ minimalne opóźnienia przetwarzania, odporność na~awarie oraz skalowalność infrastruktury.
Dokonano przeglądu istniejących rozwiązań w~obszarze przetwarzania danych finansowych, w~tym platform strumieniowych takich jak Apache Kafka czy~NATS JetStream, oraz porównano mechanizmy składowania danych w~architekturach warstwowych.
Rozdział zawiera także uzasadnienie wyboru technologii wykorzystanych w~projekcie.

Rozdział~\ref{ch:projekt} zawiera szczegółowy projekt architektury systemu, obejmujący dekompozycję na~warstwy: pozyskiwania danych, przetwarzania strumieniowego, transformacji \textit{ETL}, składowania oraz udostępniania.
Opisano strukturę baz danych MongoDB przechowujących dane operacyjne oraz~konfigurację warstwy obiektowej MinIO jako pamięci zimnej.
Przedstawiono przepływ danych między poszczególnymi komponentami, mechanizmy komunikacji asynchronicznej przez kolejki \textit{JetStream} oraz zasady wersjonowania i~retencji danych.
Rozdział uwzględnia także schemat połączenia z~interfejsem GraphQL oraz integracją z~API Telegrama w~celu dystrybucji powiadomień.

Rozdział~\ref{ch:implementacja} przedstawia szczegóły implementacyjne systemu, w~tym wybrane biblioteki języka \textit{Go} do~obsługi WebSocket, deserializacji JSON oraz interakcji z~bazami danych.
Opisano napotkane problemy, takie jak zarządzanie ponownymi połączeniami po~awariach, obsługa backpressure w~strumieniach danych czy~optymalizacja zapytań do~MongoDB przy dużym wolumenie rekordów.
Zaprezentowano także fragmenty kodu źródłowego ilustrujące kluczowe mechanizmy, w~szczególności logikę archiwizacji danych oraz~konstrukcję (\textit{resolverów}) GraphQL\@.

Rozdział~\ref{sec:rozwoj} wskazuje możliwe kierunki dalszego rozwoju systemu, obejmujące rozszerzenie o~dodatkowe giełdy kryptowalut, zastosowanie algorytmów wykrywania anomalii oraz integrację z~systemami monitoringu i~powiadamiania, takimi jak Prometheus czy~Grafana.
Rozważono także możliwość wdrożenia mechanizmów uczenia maszynowego do~predykcji ruchów cenowych oraz automatycznego generowania zleceń handlowych na~podstawie analizy danych historycznych.

Rozdział~\ref{ch:testy} dokumentuje przeprowadzone testy manualne weryfikujące poprawność działania poszczególnych komponentów.
Opisano scenariusze testowe obejmujące m.in.\ odbiór danych z~WebSocket, transformację rekordów w~procesach \textit{ETL}, obsługę awarii komponentów oraz poprawność odpowiedzi GraphQL\@.
Przedstawiono wyniki pomiarów opóźnień w~łańcuchu przetwarzania oraz statystyki dotyczące wolumenu przetwarzanych danych.

Pracę zamyka rozdział~\ref{ch:podsumowanie}, w~którym podsumowano rezultaty projektu, omówiono stopień realizacji założonych celów oraz~wskazano wnioski płynące z~realizacji systemu.
Zidentyfikowano mocne strony zaprojektowanej architektury, takie jak modularność, odporność na~awarie oraz niskie koszty operacyjne, a~także obszary wymagające dalszych udoskonaleń.

